\chapter{Asyncio Inception, Pathlib, and Enum (Python 3.4)}

Python 3.4 shipped three additions that each replaced a pile of ad-hoc idioms with a designed
abstraction: \textbf{asyncio} (PEP 3156) gave single-threaded I/O concurrency a standard event
loop; \textbf{pathlib} (PEP 428) replaced string-based path juggling with path \emph{objects};
and \textbf{enum} (PEP 435) gave Python real enumerations. For asyncio we focus on the
\textbf{event loop and the Future/Task machinery}---its 3.4 inception---and defer the
\texttt{async}/\texttt{await} \emph{syntax} to Chapter 11 and the full internals and structured
concurrency to Vol IX.

\section*{Section Index}
\begin{itemize}
  \item \textbf{10.1} Asyncio inception: the event loop, selectors, Futures, and Tasks (PEP 3156)
  \item \textbf{10.2} \texttt{pathlib}: object-oriented filesystem paths (PEP 428)
  \item \textbf{10.3} \texttt{enum}: enumerations via a metaclass (PEP 435)
  \item \textbf{10.4} Performance and anti-patterns
  \item \textbf{10.5} Summary and cross-references
\end{itemize}

\section{Asyncio inception: the event loop, selectors, Futures, and Tasks (PEP 3156)}

\textbf{Why this exists.} For I/O-bound workloads---thousands of sockets mostly \emph{waiting}---a
thread per connection wastes memory and pays context-switch and GIL-contention costs, while raw
callbacks become ``callback hell.'' \textbf{asyncio} offers a third model: \textbf{one thread, one
event loop, cooperative multitasking}. Tasks voluntarily suspend at \texttt{await} points; while
one waits on I/O, the loop runs others.

\textbf{What the loop actually does.} asyncio builds on the OS's \textbf{I/O multiplexing}: it
registers descriptors with the \texttt{selectors} module (mapping to \texttt{epoll}/\texttt{kqueue}/
\texttt{select}), calls \texttt{selector.select()} which \textbf{blocks the single thread in the
kernel} until a descriptor is ready, then runs the ready descriptors' callbacks to resume suspended
tasks. A \textbf{\texttt{Future}} holds a value that will exist later (PENDING $\to$
FINISHED/CANCELLED, with done-callbacks); a \textbf{\texttt{Task}} is a \texttt{Future} subclass
that drives a coroutine---its \texttt{\_step} calls \texttt{coro.send(...)}, registers itself as the
done-callback of the awaited future, and yields to the loop.

\begin{lstlisting}[language=Python, caption={3.4-era generator coroutine. NOT executed --- @asyncio.coroutine was REMOVED in 3.11. Shown to connect asyncio to Chapter 9's yield from.}]
import asyncio

@asyncio.coroutine
def fetch_data():
    yield from asyncio.sleep(1)     # delegate to asyncio.sleep's awaitable
    return "payload_data"
\end{lstlisting}

\begin{lstlisting}[language=Python, caption={The legacy decorator no longer exists.}]
import asyncio
print("hasattr(asyncio, 'coroutine'):", hasattr(asyncio, "coroutine"))
\end{lstlisting}

Verified output (CPython 3.13.5):

\begin{lstlisting}[caption={Output}]
hasattr(asyncio, 'coroutine'): False
\end{lstlisting}

The \textbf{modern equivalent} uses Chapter 11's \texttt{async}/\texttt{await} but the same
event-loop model. Three 0.05\,s sleeps run \textbf{concurrently}, finishing in $\sim$0.05\,s:

\begin{lstlisting}[language=Python, caption={The event loop overlaps waits --- concurrency on a single thread.}]
import asyncio, time

async def worker(name, delay):
    await asyncio.sleep(delay)
    return f"{name} done"

async def main():
    t0 = time.perf_counter()
    results = await asyncio.gather(worker("A", 0.05), worker("B", 0.05), worker("C", 0.05))
    return results, time.perf_counter() - t0

results, elapsed = asyncio.run(main())
print("results:", results)
print("elapsed (~0.05s, not 0.15s):", round(elapsed, 3), "s")
\end{lstlisting}

Verified output (CPython 3.13.5):

\begin{lstlisting}[caption={Output}]
results: ['A done', 'B done', 'C done']
elapsed (~0.05s, not 0.15s): 0.051 s
\end{lstlisting}

\textbf{The senior-engineer contrast.} This is cooperative scheduling, like Go goroutines or
Node's loop---but explicit: a task runs uninterrupted until it \texttt{await}s, so no preemption
and no lock for data touched only between awaits. The flip side (§10.4) is that a \emph{blocking}
call stalls the \textbf{entire} loop. Syntax is Chapter 11; internals, \texttt{TaskGroup}, and
structured concurrency are Vol IX.

\section{\texttt{pathlib}: object-oriented filesystem paths (PEP 428)}

\textbf{Why this exists.} Before 3.4 a path was a \texttt{str} manipulated by scattered
\texttt{os.path} functions. \textbf{pathlib} makes a path a typed object with an overloaded
\texttt{/} join and a split between \emph{pure} (no I/O) and \emph{concrete} (filesystem)
operations: \texttt{PurePath} does algebra (and parses foreign paths); \texttt{Path} adds I/O and
instantiates \texttt{PosixPath}/\texttt{WindowsPath} for the host.

\begin{lstlisting}[language=Python, caption={Path objects --- / for joining, rich accessors, cross-platform pure paths.}]
from pathlib import Path, PureWindowsPath

p = Path("/var") / "log" / "nginx.log"     # __truediv__ joins
print("joined:", p)
print("parts:", p.parts, "| name:", p.name, "| suffix:", p.suffix, "| stem:", p.stem)
print("with_suffix:", p.with_suffix(".gz"))
print("parent:", p.parent)
print("Windows path parsed on this host:", PureWindowsPath("C:/a/b"))
print("Path('.') concrete type:", type(Path(".")).__name__)
\end{lstlisting}

Verified output (CPython 3.13.5):

\begin{lstlisting}[caption={Output}]
joined: /var/log/nginx.log
parts: ('/', 'var', 'log', 'nginx.log') | name: nginx.log | suffix: .log | stem: nginx
with_suffix: /var/log/nginx.gz
parent: /var/log
Windows path parsed on this host: C:\a\b
Path('.') concrete type: PosixPath
\end{lstlisting}

The \texttt{/} operator handles separators per-OS, so the same code is correct on Windows and
POSIX. \texttt{Path} unifies \texttt{read\_text}, \texttt{write\_bytes}, \texttt{glob},
\texttt{mkdir}, \texttt{exists}, \texttt{stat}. Prefer \texttt{pathlib} over \texttt{os.path}.

\section{\texttt{enum}: enumerations via a metaclass (PEP 435)}

\textbf{Why this exists.} Before 3.4, ``enums'' were bare constants---no namespacing, type safety,
iteration, or readable \texttt{repr}. \textbf{enum} provides named, singleton, immutable, iterable
members. \textbf{What the metaclass does:} \texttt{EnumMeta} uses a custom namespace
(\texttt{\_EnumDict}) to capture members in order and reject duplicates, replaces each
\texttt{NAME = value} with a \textbf{singleton instance}, and installs lookup hooks.

\begin{lstlisting}[language=Python, caption={enum members are singletons, iterable, looked up by name or value, and immutable.}]
from enum import Enum

class Color(Enum):
    RED = 1
    GREEN = 2
    BLUE = 3

print("members:", [c.name for c in Color])          # __iter__ in declaration order
print("by name  Color['RED']:", Color["RED"])        # __getitem__
print("by value Color(2):", Color(2))                # __call__
print("singleton:", Color.RED is Color(1))           # same object every time
try:
    Color.RED.value = 99                              # immutable
except AttributeError as e:
    print("immutable ->", type(e).__name__)
\end{lstlisting}

Verified output (CPython 3.13.5):

\begin{lstlisting}[caption={Output}]
members: ['RED', 'GREEN', 'BLUE']
by name  Color['RED']: Color.RED
by value Color(2): Color.GREEN
singleton: True
immutable -> AttributeError
\end{lstlisting}

Compare members with \texttt{is}/\texttt{==} and switch on them in \texttt{match}/\texttt{case}
(Chapter 15). \texttt{IntEnum}, \texttt{Flag}/\texttt{IntFlag}, \texttt{auto()}, and
\texttt{StrEnum} (3.11) are developed in Vol XII.

\section{Performance and anti-patterns}

\begin{itemize}
  \item \textbf{Never block the event loop}---a sync call or CPU loop in a coroutine freezes every
    task. Use async libraries or offload via \texttt{run\_in\_executor}/\texttt{asyncio.to\_thread}
    (Vol IX). The most common asyncio bug.
  \item \textbf{Don't conflate paradigms}---asyncio is I/O concurrency; CPU parallelism needs
    processes or the free-threaded build (Vol VII).
  \item \textbf{\texttt{pathlib} has slight overhead vs raw \texttt{str}}---each op builds a new
    \texttt{Path}; in million-path loops \texttt{os.path} can be faster, otherwise prefer
    \texttt{pathlib}.
  \item \textbf{Enum pitfalls}---duplicate values create \textbf{aliases} (use \texttt{@enum.unique}
    to forbid); \texttt{IntEnum} members compare equal to ints (prefer plain \texttt{Enum} unless
    needed).
\end{itemize}

\section{Summary and cross-references}

\begin{itemize}
  \item \textbf{asyncio} (PEP 3156) is a \textbf{single-threaded cooperative event loop} over OS
    multiplexing; \textbf{Futures} hold pending results, \textbf{Tasks} drive coroutines. The 3.4
    \texttt{@asyncio.coroutine}+\texttt{yield from} syntax is gone (removed 3.11); the model is
    unchanged.
  \item \textbf{pathlib} (PEP 428): \texttt{PurePath} (algebra) vs \texttt{Path} (I/O), the
    \texttt{/} join, a rich API. Prefer over \texttt{os.path}.
  \item \textbf{enum} (PEP 435) builds \textbf{singleton, immutable, iterable} members via
    \texttt{EnumMeta}, with lookup by name and value.
\end{itemize}

\textbf{Cross-references.} \texttt{yield from} and the suspendable frame $\rightarrow$ Chapter 9.
Native \texttt{async}/\texttt{await} $\rightarrow$ Chapter 11. Event-loop internals,
\texttt{TaskGroup}, structured concurrency $\rightarrow$ Vol IX. The GIL $\rightarrow$ Vol VII. The
full \texttt{enum} family and \texttt{graphlib} $\rightarrow$ Vol XII. \texttt{pathlib}/\texttt{os}
reference $\rightarrow$ Vol XIV.
